\documentclass[border=10pt]{standalone}
\usepackage{fontspec}
\setmainfont{Latin Modern Roman}[Ligatures=TeXOff]
\XeTeXgenerateactualtext=1
\newfontfamily\bengalifont{Noto Serif Bengali}[Script=Bengali, NFSSFamily=bnfam]
\usepackage[Latin,Bengali]{ucharclasses}
% group-free font switching (as in header-common.tex): robust inside braces, cells and quotes
\makeatletter
\def\BN@saved{\rmdefault}
\def\BN@bn{bnfam}
\newcommand\BNin{\edef\BN@tmp{\f@family}\ifx\BN@tmp\BN@bn\else\let\BN@saved\BN@tmp\fi\fontfamily{bnfam}\selectfont}
\newcommand\BNout{\fontfamily{\BN@saved}\selectfont}
\makeatother
\setTransitionsFor{Bengali}{\BNin}{\BNout}
% Typographic quotes and dashes are not in a ucharclasses class (class 0), so a Bangla word right
% after “ or ‘ would not trigger the Bangla font. Treat them as Latin so that every switch fires.
\XeTeXcharclass"2018=\XeTeXcharclass`A
\XeTeXcharclass"2019=\XeTeXcharclass`A
\XeTeXcharclass"201C=\XeTeXcharclass`A
\XeTeXcharclass"201D=\XeTeXcharclass`A
\XeTeXcharclass"2013=\XeTeXcharclass`A
\XeTeXcharclass"2014=\XeTeXcharclass`A
\XeTeXcharclass"2026=\XeTeXcharclass`A

\usepackage[table]{xcolor}
\usepackage{booktabs,tabularx,array}
\definecolor{bapilrule}{HTML}{1F2937}
\definecolor{bapilstripe}{HTML}{F3F4F6}
\renewcommand{\arraystretch}{1.28}
\setlength{\tabcolsep}{5pt}

\begin{document}
\begin{minipage}{17cm}
{\large\bfseries Table 6.}\ {\large Strengths, weaknesses, opportunities and threats (SWOT)}\par\vspace{7pt}
\small
\rowcolors{2}{bapilstripe}{white}
\begin{tabularx}{\linewidth}{@{}>{\hsize=0.377\hsize\raggedright\arraybackslash}X>{\hsize=1.623\hsize\raggedright\arraybackslash}X@{}}
\arrayrulecolor{bapilrule}\toprule
\rowcolor{white}\textbf{Category} & \textbf{Factors} \\
\midrule
\textbf{Strengths} & Large digital service footprint (national portal, e-services, e-Nothi); experience with digital transformation at scale; a large, young and digitally active population; active Bangla NLP research community; existing national AI programmes \\
\textbf{Weaknesses} & Scattered and often scanned official documents; inconsistent metadata and licences; scarce compute; limited AI assurance skills; few independent evaluations; dependence on project funding \\
\textbf{Opportunities} & National AI policy drafting from Q4 2026; new data laws and the national data exchange; open-weight models that run locally; Bangla as a large-language opportunity; development-partner interest in digital public goods \\
\textbf{Threats} & Political and institutional change; misinformation and synthetic media; cyber attacks and outages; vendor lock-in; erosion of trust if AI tools fail visibly; talent migration \\
\bottomrule
\end{tabularx}
\par\vspace{3pt}{\scriptsize\color{gray} Source: BAPIL Technical Blueprint, Version 1.0 (2026). CC BY 4.0.}
\end{minipage}
\end{document}
